\documentclass[11pt]{article}
\usepackage[a4paper,margin=18mm]{geometry}
\usepackage{fontspec}
\setmainfont{Latin Modern Roman}
\usepackage{amsmath,amssymb,amsthm,amscd,graphicx,color}
\usepackage{xurl}
\usepackage[unicode,hidelinks]{hyperref}
\allowdisplaybreaks
\setlength\emergencystretch{3em}

\makeatletter
\newtheorem{theorem}{Theorem}
\newtheorem*{theorem*}{Theorem}
\newtheorem{lemma}{Lemma}
\newtheorem*{lemma*}{Lemma}
\newtheorem{corollary}{Corollary}
\newtheorem*{corollary*}{Corollary}
\newtheorem{proposition}{Proposition}
\newtheorem*{proposition*}{Proposition}
\newtheorem{conjecture}{Conjecture}
\newtheorem*{conjecture*}{Conjecture}
{\theoremstyle{definition} \newtheorem{definition}{Definition}
\newtheorem*{definition*}{Definition}
\newtheorem{example}{Example}
\newtheorem*{example*}{Example}
\newtheorem{remark}{Remark}
\newtheorem*{remark*}{Remark}
\newtheorem{note}{Note}
\newtheorem*{note*}{Note}
}
\makeatother
\usepackage[mathscr]{eucal}

\makeatletter\newlabel{appendixA}{{A}{}{Original source locator}{}{}}\makeatother
\begin{document}
\begin{center}\Large Explicit unitary spectral transform of the specified bounded radial Laplace-Beltrami operator on quantum complex hyperbolic space H n,m for n and m at least two and zero less than q less than one\end{center}
\noindent\textbf{Stable ID:} 2646\par
\noindent\textbf{Authors:} Olga Bershtein; Yevgen Kolisnyk\par
\noindent\textbf{Original work:} Harmonic Analysis on Quantum Complex Hyperbolic Spaces\par
\noindent\textbf{Version:} Official SIGMA 7 (2011) article 078, DOI 10.3842/SIGMA.2011.078; journal PDF and native TeX acquired 2026-10-01, exact bytes pinned by SHA256\par
\noindent\textbf{Selected task scope:} Theorem3 restricted to positive integers n,m>=\allowbreak{}2,N=\allowbreak{}n+\allowbreak{}m,and0<q<1. The radial Laplace-Beltrami operator square\textasciicircum{}(0) defined by the original invariant pairing and Lemma5 difference formula extends bounded self-adjointly to the completion L2(dnu\_q\textasciicircum{}(0)) of finitely supported functions on q\textasciicircum{}(-2Z\_+\allowbreak{}). The original kernel transform U f(lambda)=\allowbreak{}integral\_1\textasciicircum{}infinity f(x)Phi\_l(x)rho(x)d\_(q\textasciicircum{}-2)x is unitary to L2(dsigma) and conjugates square\textasciicircum{}(0) to multiplication by lambda, with the original lambda/l relation, rho and continuous/discrete Al-Salam-Chihara spectral measure retained.\par
\noindent\textbf{Difficulty:} H2: The radial operator must first be extracted from invariant quantum differential calculus and a Hermitian pairing on noncommuting coordinate functions. Its weighted Jackson lattice produces a bounded Jacobi operator whose coefficients and normalization differ from ordinary radial Laplacians. The authors construct generalized eigenfunctions, match the full normalized Al-Salam-Chihara recurrence, and identify the continuous-plus-discrete spectral transform on the precise radial Hilbert completion.\par
\noindent\textbf{Editorial boundary note (not author text):} The full original unlabeled Theorem3 is retained for n,m at least two on the precise radial weighted finite-function completion. All quantum coordinate/invariant integration setup, differential pairing and radial Jackson weight, explicit difference coefficients, boundedness, eigenfunction and normalized Al-Salam-Chihara recurrence/kernel calculations are bundled. Named quantum calculus/invariant integral and polynomial orthogonality remain cited prerequisites. Source bounded compact-perturbation/direct-evaluation details are preserved; no whole nonradial spectral theorem or asymptotic Harish-Chandra result is selected.\par
\noindent\textbf{Source:} \url{https://sigma-journal.com/2011/078/}\quad DOI: \url{https://doi.org/10.3842/SIGMA.2011.078}\par
\noindent\textbf{License and attribution:} Copyright retained by the named authors. Published by SIGMA under CC BY 4.0: \url{https://creativecommons.org/licenses/by/4.0/}. Source: \url{https://sigma-journal.com/about.html}. Adaptation consists of standalone excerpt formatting, cover and numbering restoration; original author language and mathematical text are retained.\par
\noindent\textbf{Editorial symbol notice (not author text):} The unit is the radial operator and its specified weighted finite-function completion, not a spectral theorem for the full nonradial quantum Laplacian. Source square/overbar and q-Jackson normalization formulas remain verbatim. Native primary theorem has no label; its exact theorem environment is at lines1004-1014 and author proof extends to1113.\par
\noindent\textbf{External original locator appendixA:} Appendix A calculates Harish-Chandra asymptotics. This later interpretation is not needed for the radial spectral transform, whose complete Jacobi recurrence and orthogonality normalization are retained.\par
\medskip\hrule\medskip\noindent\textbf{Author text begins below.}\par
\setcounter{section}{1}
\setcounter{subsection}{0}
\setcounter{subsubsection}{0}
\setcounter{equation}{0}
% Original source lines 83--1113
\section{Preliminaries on quantum group theory}\label{sec_1}

Everywhere in the sequel we suppose $q\in(0,1)$. All algebras are associative and unital.

The Hopf algebra $U_q\mathfrak{sl}_{N}$ is given by its generators $K_i$, $K_i^{-1}$, $E_i$, $F_i$,
$i=1,2,\ldots,N-1$, and the relations:
\begin{gather*}
K_iK_j=K_jK_i,\qquad K_iK_i^{-1}=K_i^{-1}K_i=1,
\\
K_iE_i=q^2E_iK_i,\qquad K_iF_i=q^{-2}F_iK_i,
\\
K_iE_j=q^{-1}E_jK_i,\qquad K_iF_j=q F_jK_i, \qquad |i-j|=1,
\\
E_iF_j-F_jE_i=\delta_{ij}\frac{K_i-K_i^{-1}}{q-q^{-1}},
\\
E_i^2E_j-\big(q+q^{-1}\big)E_iE_jE_i+E_jE_i^2=0,\qquad|i-j|=1,
\\
F_i^2F_j-\big(q+q^{-1}\big)F_iF_jF_i+F_jF_i^2=0,\qquad|i-j|=1,
\\
[E_i,E_j]=[F_i,F_j]=0,\qquad|i-j|\ne 1.
\end{gather*}
 The comultiplication $\Delta$, the antipode $S$, and the counit
$\varepsilon$ are def\/ined on the generators by
\begin{gather*}
\Delta(E_i)=E_i\otimes 1+K_i\otimes E_i,\qquad \Delta(F_i)=F_i\otimes K_i^{-1}+1\otimes F_i,\qquad
\Delta(K_i)=K_i\otimes K_i,
\\
S(E_i)=-K_i^{-1}E_i,\qquad S(F_i)=-F_iK_i,\qquad S(K_i)=K_i^{-1},
\\
\varepsilon(E_i)=\varepsilon(F_i)=0,\qquad\varepsilon(K_i)=1,
\end{gather*} see \cite[Chapter 4]{Jant}.

We need also the Hopf algebra $\mathbb{C}[SL_{N}]_q$ of matrix elements of f\/inite dimensional weight
$U_q\mathfrak{sl}_{N}$-modules. Recall that $\mathbb{C}[SL_{N}]_q$ can be def\/ined by the generators $t_{ij}$,
$i,j=1,\dots,N$ (the matrix elements of the vector representation in a weight basis) and the relations
\begin{alignat*}{3}
&t_{ij'}t_{ij''}=qt_{ij''}t_{ij'},\qquad && j'<j'',&
\\ &t_{i'j}t_{i''j}=qt_{i''j}t_{i'j},\qquad && i'<i'',&
\\ &t_{ij}t_{i'j'}=t_{i'j'}t_{ij},\qquad && i<i'\;\&\;j>j', &
\\ &t_{ij}t_{i'j'}=t_{i'j'}t_{ij}+\big(q-q^{-1}\big)t_{ij'}t_{i'j},\qquad
&& i<i'\;\&\; j<j',&
\end{alignat*}
together with one more relation
\[
\det\nolimits_q\mathbf{t}=1,
\]
where $\det\nolimits_q\mathbf{t}$ is a $q$-determinant of the matrix $\mathbf{t}=(t_{ij})_{i,j=1,\dots,N}$:
\[
\det\nolimits_q\mathbf{t}=\sum\limits_{s\in S_{N}}(-q)^{l(s)}t_{1 s(1)} t_{2s(2)}\cdots t_{Ns(N)},
\]
with $l(s)=\mathrm{card}\{(i,j)|i<j\;\&\;s(i)>s(j)\}$. The algebra $\mathbb C[SL_N]_q$ is endowed with the
standard structure of $U_q^{\rm op}\mathfrak{sl}_N \otimes U_q \mathfrak{sl}_N$-module algebra (here 'op'
ref\/lects the fact that we should change the multiplication by the opposite one).

Let also $U_q\mathfrak{su}_{n,m}$, $m+n=N$, denotes the Hopf $*$-algebra $(U_q \mathfrak{sl}_N,*)$ given by
\[
(K_j^{\pm 1})^*=K_j^{\pm 1},\qquad E_j^*=
\begin{cases}
K_jF_j,& j \ne n,
\\ -K_jF_j,& j=n,
\end{cases}\qquad F_j^*=
\begin{cases}
E_jK_j^{-1},& j \ne n,
\\ -E_jK_j^{-1},& j=n,
\end{cases}
\]
with $j=1,\ldots,N-1$ \cite{RTF, polmat}.

Recall the notion of an algebra of `regular functions on the quantum principal homogeneous space' $X$
constructed in \cite{polmat}. Put $\operatorname{Pol}(\widetilde{X})_q\overset{\mathrm{def}}{=}
(\mathbb{C}[SL_N]_q,*)$, where the involution $*$ is def\/ined by
\begin{gather}\label{def_inv}
t_{ij}^*=\mathrm{sign}[(i-m-1/2)(n-j+1/2)](-q)^{j-i}\det\nolimits_qT_{ij}.
\end{gather}
Here $\det_q$ is the quantum determinant \cite{Ch-P}, and the matrix $T_{ij}$ is derived from the matrix
$\mathbf{t}=(t_{kl})$ by discarding its $i$'s row and $j$'s column. In \cite{polmat} it is proved that
$\operatorname{Pol}(\widetilde{X})_q$ is a $U_{q}\mathfrak{su}_{n,m}$-module $*$-algebra.

\section[A $*$-algebra $\operatorname{Pol}(\mathscr{H}_{n,m})_{q}$]{A $\boldsymbol{*}$-algebra $\boldsymbol{\operatorname{Pol}(\mathscr{H}_{n,m})_{q}}$}\label{sec_2}

Let $m,n\in\mathbb{N}$, $m\ge 2$, and $N\overset{\mathrm{def}}{=}n+m$. Recall that the classical complex
hyperbolic space $\mathscr{H}_{n,m}$ can be obtained by projectivization of the domain
\[
\widehat{\mathscr{H}}_{n,m}=\left\{(t_1,\ldots,t_N)\in\mathbb{C}^N\; \Big|\;
-\sum^n_{j=1}|t_j|^2+\sum^N_{j=n+1}|t_{j}|^2>0 \right\}.
\]

Now we pass from the classical case $q=1$ to the quantum case $0<q<1$. Let us consider the well known
\cite{RTF} $q$-analog of the polynomial algebras. Let $\operatorname{Pol}(\widehat{\mathscr{H}}_{n,m})_q$
be the unital $*$-algebra with the generators $t_1,t_{2},\ldots,t_N$ and the commutation relations as
follows:
\begin{gather*}
t_it_j = qt_jt_i,\qquad i<j,
\\ t_it_j^* = qt_j^*t_i,\qquad i\ne j,
\\ t_it_i^* = t_i^*t_i+\big(q^{-2}-1\big)\sum_{k=i+1}^Nt_kt_k^*,\qquad i>n,
\\ t_it_i^* = t_i^*t_i+\big(q^{-2}-1\big)\sum_{k=i+1}^nt_kt_k^*-
\big(q^{-2}-1\big)\sum_{k=n+1}^Nt_kt_k^*,\qquad i\le n.
\end{gather*}
Obviously, \[c=-\sum_{j=1}^nt_jt_j^*+\sum_{j=n+1}^Nt_jt_j^*\] is central in
$\operatorname{Pol}(\widehat{\mathscr{H}}_{n,m})_q$. Moreover, $c$ is not a zero divisor in
$\operatorname{Pol}(\widehat{\mathscr{H}}_{n,m})_q$. This allows one to embed the $*$-algebra
$\operatorname{Pol}(\widehat{\mathscr{H}}_{n,m})_q$ into its localization
$\operatorname{Pol}(\widehat{\mathscr{H}}_{n,m})_{q,c}$ with respect to the multiplicative system
$c^{\mathbb{N}}$.

The $*$-algebra $\operatorname{Pol}(\widehat{\mathscr{H}}_{n,m})_{q,c}$ admits the following bigrading:
\[\deg t_{j}=(1,0),\qquad\deg t_j^*=(0,1),\qquad j=1,2,\ldots,N.\] Introduce the notation
\[
\operatorname{Pol}(\mathscr{H}_{n,m})_q=\left\{ f\in
\operatorname{Pol}(\widehat{\mathscr{H}}_{n,m})_{q,c} \; \big|\; \deg f=(0,0)\right\}.
\]
This $*$-algebra $\operatorname{Pol}(\mathscr{H}_{n,m})_q$ will be called the algebra of regular functions on
the quantum hyperbolic space.



We are going to endow the $*$-algebra $\operatorname{Pol}(\mathscr{H}_{n,m})_q$ with a structure of
$U_{q}\mathfrak{su}_{n,m}$-module algebra~\cite{Ch-P}. For this purpose, we embed it into the
$U_{q}\mathfrak{su}_{n,m}$-module $*$-algebra $\operatorname{Pol}(\widetilde{X})_q$.

By a $q$-analog of the Laplace expansion of $\det_q \mathbf{t}$ along the f\/irst row \cite[Section~9.2]{Kl-Sch} and
\eqref{def_inv}, one can obtain from $\det_q \mathbf{t}=1$ that
\[
-\sum_{j=1}^nt_{1j}t_{1j}^*+\sum_{j=n+1}^Nt_{1j}t_{1j}^*=1.
\]
Thus the map $J:t_j\mapsto t_{1j}$, $j=1,2,\ldots,N$, admits a unique extension to a homomorphism of
$*$-algebras $J:\operatorname{Pol}(\widehat{\mathscr{H}}_{n,m})_{q,c}\to
\operatorname{Pol}(\widetilde{X})_q$. Its image will be denoted by
$\operatorname{Pol}(\widetilde{\mathscr{H}}_{n,m})_q$. It is easy to verify that the $*$-algebra
$\operatorname{Pol}(\mathscr{H}_{n,m})_q$ is {\it embedded} this way into
$\operatorname{Pol}(\widetilde{\mathscr{H}}_{n,m})_q$ and its image is just the subalgebra in
$\operatorname{Pol}(\widetilde{\mathscr{H}}_{n,m})_q$ generated by $t_{1j}t_{1k}^*$, $j,k=1,2,\ldots,N$
(recall that $c$ goes to $\det_q \mathbf{t}=1$). In what follows we will identify
$\operatorname{Pol}(\mathscr{H}_{n,m})_q$ with its image under the map $J$.

Consider the subalgebra $U_q \mathfrak{s(gl}_1 \times \mathfrak{gl}_{N-1})$ generated by $K_i^{\pm 1}$,
$i=1,\dots,N-1,$ $E_j$, $F_j$, $j =2,\dots,N-1.$ By obvious reasons,
\[
\operatorname{Pol}(\mathscr{H}_{n,m})_q=\left\{f \in \operatorname{Pol}(\widetilde{X})_q \;\big|\;
L(\xi)f=\varepsilon(\xi)f, \xi \in U_q \mathfrak{s(gl}_1 \times \mathfrak{gl}_{N-1})\right\},
\]
where $L$ is the left action of $U_q^{\rm op} \mathfrak{sl}_N$ in $\operatorname{Pol}(\widetilde{X})_q$.

Let $I_\varphi$, $\varphi\in\mathbb{R}/2\pi\mathbb{Z}$, be the $*$-automorphism of the $*$-algebra
$\operatorname{Pol}(\widetilde{\mathscr{H}}_{n,m})_q$ def\/ined on the generators $\{t_j\}_{j=1,\ldots,N}$ by
\begin{gather}\label{I_phi}
I_\varphi: \ t_j\mapsto e^{i\varphi}t_j.
\end{gather}
We
use the notation $t_j$ instead of $t_{1j}$ for the generators of
$\operatorname{Pol}(\widetilde{\mathscr{H}}_{n,m})_q$.

Then one more description of $\operatorname{Pol}(\mathscr{H}_{n,m})_q$ is as follows:
\[
\operatorname{Pol}(\widetilde{\mathscr{H}}_{n,m})_q = \left\{ f\in
\operatorname{Pol}(\widetilde{\mathscr{H}}_{n,m})_q\; \big| \;I_\varphi(f)=f\text{ \ for all \
}\varphi\right\}.
\]

At the end of this section we list explicit formulas for the action of $U_{q}\mathfrak{su}_{n,m}$ on
$\operatorname{Pol}(\widetilde{\mathscr{H}}_{n,m})$:
\begin{gather*}
E_jt_i =
\begin{cases}
q^{-1/2}t_{i-1}, & j+1=i,
\\ 0, & \text{otherwise},
\end{cases}\qquad
F_jt_i =
\begin{cases}
q^{1/2}t_{i+1}, & j=i,
\\ 0, & \text{otherwise},
\end{cases}
\\ K_j^{\pm 1}t_i  =
\begin{cases}
q^{\pm 1}t_i, & j=i,
\\ q^{\mp 1}t_i, & j+1=i,
\\ t_i, & \text{otherwise}.
\end{cases}
\end{gather*}



\section[Algebras of generalized and finite functions on the quantum $\mathscr{H}_{n,m}$]{Algebras of generalized and f\/inite functions\\ on the quantum $\boldsymbol{\mathscr{H}_{n,m}}$}\label{finite}

Let us construct a faithful $*$-representation $T$ of $\operatorname{Pol}(\mathscr{H}_{n,m})_q$ in a
pre-Hilbert space $\mathscr{H}$ (our method is well known; see, for example, \cite{polmat}).

The space $\mathscr{H}$ is a linear span of its orthonormal basis
$\{e(i_1,i_2,\ldots,i_{N-1})\, |\,i_1,\ldots,i_n\in-\mathbb{Z}_+$; $i_{n+1}, \ldots,i_{N-1}\in\mathbb{N}\}$.

The $*$-representation $T$ is a restriction to $\operatorname{Pol}(\mathscr{H}_{n,m})_q$ of the
$*$-representation of $\operatorname{Pol}(\widetilde{\mathscr{H}}_{n,m})_q$ def\/ined by
\begin{gather*}
T(t_j)e(i_1,\ldots,i_{N-1})  = q^{\sum\limits_{k=1}^{j-1}i_k} \big(q^{2(i_j-1)}-1\big)^{1/2}
e(i_1,\ldots,i_j-1,\ldots,i_{N-1}),
\\ T(t_j^*)e(i_1,\ldots,i_{N-1})  =
q^{\sum\limits_{k=1}^{j-1}i_k} \big(q^{2i_j}-1\big)^{1/2} e(i_1,\ldots,i_j+1,\ldots,i_{N-1}),
\end{gather*}
for $j\le n$,
\begin{gather*}
T(t_j)e(i_1,\ldots,i_{N-1})  = q^{\sum\limits_{k=1}^{j-1}i_k} \big(1-q^{2(i_j-1)}\big)^{1/2}
e(i_1,\ldots,i_j-1,\ldots,i_{N-1}),
\\ T(t_j^*)e(i_1,\ldots,i_{N-1})  =
q^{\sum\limits_{k=1}^{j-1}i_k} \big(1-q^{2i_j}\big)^{1/2} e(i_1,\ldots,i_j+1,\ldots,i_{N-1}),
\end{gather*}
for $n<j<N$, and, f\/inally,
\begin{gather*}
T(t_N)e(i_1,\ldots,i_{N-1})  = q^{\sum\limits_{k=1}^{N-1}i_k}e(i_1,\ldots,i_{N-1}),
\\ T(t_N^*)e(i_1,\ldots,i_{N-1})  =
q^{\sum\limits_{k=1}^{N-1}i_k}e(i_1,\ldots,i_{N-1}).
\end{gather*}

Def\/ine the elements $\{x_j\}_{j=1,\ldots,N}$ as follows:
\begin{gather}\label{x_j}
x_j\overset{\mathrm{def}}{=}
\begin{cases}
\sum\limits_{k=j}^Nt_kt_k^*, & j>n,
\\ -\sum\limits_{k=j}^nt_kt_k^*+\sum\limits_{k=n+1}^Nt_kt_k^*, & j\le n.
\end{cases}
\end{gather}
Obviously, $x_1=1$, $x_ix_j=x_jx_i$,
\begin{gather}\label{t_jx_k}
t_jx_k=
\begin{cases}
q^2x_kt_j, & j<k,
\\ x_kt_j, & j\ge k.
\end{cases}
\end{gather}

The vectors $e(i_1,\ldots,i_{N-1})$ are joint eigenvectors of the operators $T(x_j)$, $j=1,2,\ldots,N$:
\begin{gather*}%\label{Tx_j}
T(x_j)e(i_1,\ldots,i_{N-1})= q^{2\sum\limits_{k=1}^{j-1}i_k}e(i_1,\ldots,i_{N-1}).
\end{gather*}

The joint spectrum of the pairwise commuting operators $T(x_j)$, $j=1,2,\ldots,N$, is
\begin{gather*}
\mathfrak{M}=\big\{(x_1,\ldots,x_N)\in\mathbb{R}^N \;| \; x_i/x_j\in q^{2\mathbb{Z}}\;\&\;1=x_1\le x_2\le\cdots\le
x_{n+1}, \\
\phantom{\mathfrak{M}=\big\{}{} \&\; x_{n+1}>x_{n+2}>\cdots>x_N>0\big\}.
\end{gather*}

The next proposition was proved in \cite{BS}.

\begin{proposition}
$T$ is a faithful representation of $\operatorname{Pol}(\mathscr{H}_{n,m})_q$.
\end{proposition}



Let us now introduce the notion of generalized functions on the quantum complex hyperbolic space
$\mathscr{H}_{n,m}$. Evidently, using the commutation relations, one can decompose every polynomial $f \in
\operatorname{Pol}(\widetilde{\mathscr{H}}_{n,m})_q$ as follows:
\begin{gather}\label{decomp_pol_Hnm_0}
f= \!\!\!\sum_{I=(i_1,\ldots,i_N), J=(j_1,\ldots,j_N) \in \mathbb Z_+^N}\!\!\! c_{IJ}t_1^{i_1}\cdots
t_n^{i_n}t_{n+1}^{*i_{n+1}}\cdots t_N^{*i_N}t_N^{j_N}\cdots t_{n+1}^{j_{n+1}}t_n^{*j_n}\cdots t_1^{*j_1},
\!\!\!\qquad c_{IJ} \in \mathbb C.\!\!\!
\end{gather}
Due to \eqref{x_j}, the latter can be reduced to the decomposition{\samepage
\begin{gather}\label{decomp_pol_Hnm}
f=\sum_{(i_1,\ldots,i_N,j_1,\ldots,j_N):\;i_kj_k=0} t_1^{i_1}\cdots t_n^{i_n}t_{n+1}^{*i_{n+1}}\cdots
t_N^{*i_N}f_{IJ}(x_2,\ldots,x_N)t_N^{j_N}\cdots t_{n+1}^{j_{n+1}}t_n^{*j_n}\cdots t_1^{*j_1},
\end{gather}
where $f_{IJ}(x_2,\ldots,x_N)$ are polynomials.}

One can equip $\operatorname{Pol}(\widetilde{\mathscr{H}}_{n,m})_q$ with the weakest topology such that the
functionals
\[
l_{(i_1,\ldots,i_N;j_1,\ldots,j_N)}(f)=(T(f)e(i_1,\ldots,i_N),e(j_1,\ldots,j_N))
\]
are continuous. The completion of $\operatorname{Pol}(\widetilde{\mathscr{H}}_{n,m})_q$ w.r.t.\ this topology
will be considered as the space of generalized functions on the quantum $\widehat{\mathscr{H}}_{n,m}$ and
denoted by $\mathscr{D}(\widetilde{\mathscr{H}}_{n,m})_q'$. Naturally, one can extend $T$ to a representation
of $\mathscr{D}(\widetilde{\mathscr{H}}_{n,m})_q'$ by continuity. Now \eqref{decomp_pol_Hnm} allows one to
identify $\mathscr{D}(\widetilde{\mathscr{H}}_{n,m})_q'$ with the space of formal series
\[
f=\sum_{(i_1,\ldots,i_N,j_1,\ldots,j_N):\;i_kj_k=0} t_1^{i_1}\cdots
t_n^{i_n}t_{n+1}^{*i_{n+1}}\cdots t_N^{*i_N}f_{IJ}(x_2,\ldots,x_N)t_N^{j_N}\cdots
t_{n+1}^{j_{n+1}}t_n^{*j_n}\cdots t_1^{*j_1},
\]
where $f_{IJ}(x_2,\ldots,x_N)$ are functions on $\mathfrak{M}$. The topology on this space of formal series
is the topology of pointwise convergence of the functions $f_{IJ}$.


Denote by $f_0$ the following function
\begin{gather}\label{f_0(x_n+1)}
f_0=f_0(x_{n+1})=
\begin{cases}
1, & x_{n+1}=1,
\\ 0, & x_{n+1}\in q^{-2\mathbb{N}}.
\end{cases}
\end{gather}
(Recall that $\mathrm{spec}\,x_{n+1}=q^{-2\mathbb{Z}_+}$.) Thus $f_0$ is a $q$-analog of the characteristic
function of the submanifold
\[
\big\{ (t_1,\ldots,t_N)\in\mathbb{C}^N\; \big|\;t_1=t_2=\cdots=t_n=0 \big\}\cap\mathscr{H}_{n,m}.
\]

Introduce now a $*$-algebra $\operatorname{Fun}(\widetilde{\mathscr{H}}_{n,m})_q \subset
\mathscr{D}(\widetilde{\mathscr{H}}_{n,m})_q'$ generated by
$\operatorname{Pol}(\widetilde{\mathscr{H}}_{n,m})_q$ and $f_0$. Easy computations from~\eqref{t_jx_k} show
that $f_0$ satisf\/ies the following relations:
\begin{gather*}
 t_j^*f_0=f_0t_j=0,\qquad j\le n,
\\  x_{n+1}f_0=f_0x_{n+1}=f_0,
\\  f_0^2=f_0^*=f_0,
\\  t_jf_0=f_0t_j,\qquad t_j^*f_0=f_0t_j^*,\qquad j\ge n+1.
\end{gather*}

The relation $I_\varphi f_0=f_0$ allows one to extend the $*$-automorphism $I_\varphi$ \eqref{I_phi} of the
algebra $\operatorname{Pol}(\widetilde{\mathscr{H}}_{n,m})_q$ to the $*$-automorphism of
$\operatorname{Fun}(\widetilde{\mathscr{H}}_{n,m})_q$. Let
\[
\operatorname{Fun}(\mathscr{H}_{n,m})_q\overset{\mathrm{def}}{=} \left\{ f\in
\operatorname{Fun}(\widetilde{\mathscr{H}}_{n,m})_q\;\big|\; I_\varphi f=f\right\}.
\]
Obviously, there exists a unique extension of the $*$-representation $T$ to a $*$-representation of the
$*$-algebra $\operatorname{Fun}(\mathscr{H}_{n,m})_q$ such that $T(f_0)$ is the orthogonal projection of
$\mathscr{H}$ onto the linear span of vectors $\{e(\underbrace{0,\ldots,0}_n,i_{n+1},\ldots,i_{N-1})|\:
i_{n+1},\ldots,i_{N-1}\in\mathbb{N}\}$.



Let $\mathscr{D}(\mathscr{H}_{n,m})_{q,\mathfrak{k}}$ be the two-sided ideal of
$\operatorname{Fun}(\mathscr{H}_{n,m})_q$ generated by $f_0$. We call this ideal the algebra of f\/inite
functions on the quantum hyperbolic space. It is a quantum analog for the algebra of $U \mathfrak{k}=U
\mathfrak{s(gl}_n \times \mathfrak{gl}_m)$-f\/inite smooth functions on $\mathscr{H}_{n,m}$ with compact
support.

\begin{remark}
Let us explain the adjective `f\/inite'. If $f$ is a f\/inite function, $T(f)$ is an operator with only a f\/inite
number of nonzero entries. However, we do not consider all possible f\/inite functions (and, therefore, all
operators with f\/inite number of nonzero entries) but only $U_q \mathfrak{k}$-f\/inite ones, cf.~\cite{Faraut}.
\end{remark}

It was proved in \cite{BS} that

\begin{theorem}
The representation $T$ of $\mathscr{D}(\mathscr{H}_{n,m})_{q,\mathfrak{k}}$ is faithful.
\end{theorem}

\begin{remark}\label{remark2}
Let $f(x_{n+1})$ be a polynomial. Then it follows from \eqref{x_j}, \eqref{t_jx_k} that
\begin{gather}\label{t_ift_^*}
\sum_{i=1}^nt_if(x_{n+1})t_i^*=f\big(q^2x_{n+1}\big)\sum_{i=1}^nt_it_i^*= f\big(q^2x_{n+1}\big)(x_{n+1}-1).
\end{gather}
This computation, together with \eqref{f_0(x_n+1)}, allows one to consider the element
$f_1=\sum\limits_{i=1}^nt_if_0t_i^*$ as a~function of $x_{n+1}$ such that
\[
f_1(x_{n+1})=
\begin{cases}
q^{-2}-1, & x_{n+1}=q^{-2},
\\ 0, & x_{n+1}=1\text{ \ or \ }x_{n+1}\in q^{-2\mathbb{N}-2}.
\end{cases}
\]

Thus a multiple application of \eqref{t_ift_^*} leads to the following claim:
$\mathscr{D}(\mathscr{H}_{n,m})_{q,\mathfrak{k}}$ contains {\it all} f\/inite functions of $x_{n+1}$ (i.e.,
such functions $f$ that $f(q^{-n})=0$ for all but f\/initely many $n\in\mathbb{N}$).
\end{remark}

The $U_{q}\mathfrak{su}_{n,m}$-module algebra structure is established on
$\mathscr{D}(\mathscr{H}_{n,m})_{q,\mathfrak{k}}$ by the following:
\begin{gather*}
E_nf_0 = -\frac{q^{-1/2}}{q^{-2}-1}t_nf_0t_{n+1}^*, \qquad F_nf_0 = -\frac{q^{3/2}}{q^{-2}-1}t_{n+1}f_0t_n^*,
\qquad K_nf_0 = f_0,
\\
E_jf_0 = F_jf_0=(K_j-1)f_0=0,\qquad j\ne n.
\end{gather*}



Now we present an explicit formula for a positive invariant integral on the space of f\/inite functions
$\mathscr{D}(\mathscr{H}_{n,m})_{q,\mathfrak{k}}$ and thereby establish its existence.

Let $\nu_q:\mathscr{D}(\mathscr{H}_{n,m})_{q,\mathfrak{k}} \to\mathbb{C}$ be a linear functional def\/ined by
\[
\nu_q(f)=\operatorname{Tr}(T(f)\cdot Q)=\int_{\mathscr{H}_{n,m}} fd\nu_q,
\]
where $Q:\mathscr{H}\to\mathscr{H}$ is the linear operator given on the basis elements
$e(i_1,\ldots,i_{N-1})$ by
\[
Qe(i_1,\ldots,i_{N-1})= \mathrm{const_1} \, q^{2\sum\limits_{j=1}^{N-1}(N-j)i_j}e(i_1,\ldots,i_{N-1}),
\qquad \mathrm{const_1}>0.
\]

\begin{theorem}[\protect{\cite{BS}}]
The functional $\nu_q$ is well defined, positive, and $U_{q}\mathfrak{su}_{n,m}$-invariant.
\end{theorem}

One has to normalize this integral in a some way. In~\cite{BS} we put $\int_{\mathscr{H}_{n,m}} f_0
d\nu_q=1$, so the constant in the previous theorem equals
\[
\mathrm{const_1}= \prod\limits_{j=1}^{m-1}\big(q^{-2j}-1\big).
\]



\section[The subalgebra $\mathscr{D}(\mathscr{H}_{n,m})_{q,\mathfrak{k}}^{U_q \mathfrak{k}}$]{The subalgebra $\boldsymbol{\mathscr{D}(\mathscr{H}_{n,m})_{q,\mathfrak{k}}^{U_q \mathfrak{k}}}$}\label{sec_4}

In this section we restrict ourselves to subalgebras of $U_q \mathfrak k$-invariant elements. It is well
known that $\mathscr{H}_{n,m}$ is a pseudo-Hermitian symmetric space of rank~1~\cite{Molch2}. The following
proposition is a~natural quantum analog for this fact.

\begin{proposition}
$\mathscr{D}(\mathscr{H}_{n,m})_{q,\mathfrak{k}}^{U_q \mathfrak{k}}=\{f(x_{n+1}) \in
\mathscr{D}(\mathscr{H}_{n,m})_{q,\mathfrak{k}}\}$.
\end{proposition}
Since $U_q \mathfrak{k}=\mathbb C[K_n,K_n^{-1}] \otimes U_q \mathfrak{sl}_n \otimes U_q \mathfrak{sl}_m$, the
proof of this proposition follows from the next statement on $U_q \mathfrak{sl}_n \times U_q
\mathfrak{sl}_m$-isotypic components in $\mathscr{D}(\mathscr{H}_{n,m})_{q,\mathfrak{k}}$.
\begin{proposition}
$U_q \mathfrak{sl}_n \times U_q \mathfrak{sl}_m$-isotypic components of
$\mathscr{D}(\mathscr{H}_{n,m})_{q,\mathfrak{k}}$ correspond to the modules $L^{(n)}(a\varpi_1+d\varpi_{n-1})
\boxtimes L^{(m)}(c\varpi_1 +b \varpi_{m-1})$ with $a,b,c,d \in \mathbb Z_+$, $a+c=b+d$ $($with infinite
multiplicity$)$. The highest weight subspace is spanned by the vectors
$t_1^at_N^{*b}\varphi(x_{n+1})t_{n+1}^ct_n^{*d}$, where $\varphi(x_{n+1})$ is a finite function.
\end{proposition}

\begin{remark}
The sign $\boxtimes$ ref\/lects the fact that the multipliers are modules of dif\/ferent algebras $U_q
\mathfrak{sl}_n$ and $U_q \mathfrak{sl}_m$.
\end{remark}

\begin{proof} Let us describe $U_q \mathfrak{sl}_n \times U_q \mathfrak{sl}_m$-highest weight vectors in
$\mathscr{D}(\mathscr{H}_{n,m})_{q,\mathfrak{k}}$. One can decompose every f\/inite function $f \in
\mathscr{D}(\mathscr{H}_{n,m})_{q,\mathfrak{k}}$ in the following way
\begin{gather*}
f= \sum_{\substack{I=(i_1,\ldots,i_N), J=(j_1,\ldots,j_N) \in \mathbb Z_+^N \\ i_1+\cdots +
i_n+j_{n+1}+\cdots +j_N\\ = i_{n+1}+\cdots+i_N+j_1+\cdots +j_n}} \!\!\!\!\!c_{IJ}t_1^{i_1}\cdots
t_n^{i_n}t_{n+1}^{*i_{n+1}}\cdots t_N^{*i_N}f_0t_N^{j_N}\cdots t_{n+1}^{j_{n+1}}t_n^{*j_n}\cdots t_1^{*j_1},
\qquad c_{IJ} \in \mathbb C.
\end{gather*}
(just by applying the commutation relations, cf.~\eqref{decomp_pol_Hnm_0}). Let $L^{(n)}(\lambda)$ be the
f\/inite dimensional $U_q \mathfrak{sl}_n$-module with highest weight $\lambda$ and $\varpi_j$ are fundamental
weights of $\mathfrak{sl}_n$. By standard arguments,
\begin{gather*}
 {\rm l.s.}\big\{t_1^{l_1}t_2^{l_2}\cdots t_n^{l_n}\,|\,l_1+\cdots+l_n=a\big\}=L^{(n)}(a\varpi_1),
\\  {\rm l.s.}\big\{t_1^{*l_1}t_2^{*l_2}\cdots t_n^{*l_n}\,|\,l_1+\cdots+l_n=a\big\}=L^{(n)}(a\varpi_{n-1}),
\\  {\rm l.s.}\big\{t_{n+1}^{l_{n+1}}t_{n+2}^{l_{n+2}}\cdots t_N^{l_N}\,|\, l_{n+1}+\cdots+l_N=a\big\}=L^{(m)}(a\varpi_1),
\\  {\rm l.s.}\big\{t_{n+1}^{*l_{n+1}}t_{n+2}^{*l_{n+2}}\cdots
t_N^{*l_N}\,|\,l_{n+1}+\cdots+l_N=a\big\}=L^{(m)}(a\varpi_{m-1}).
\end{gather*}
Thus we have an epimorphism
\begin{gather*}
\bigoplus_{a+c=b+d} L^{(n)}(a\varpi_1) \otimes  L^{(m)}(d\varpi_{m-1}) \otimes L^{(m)}(c\varpi_1) \otimes
L^{(n)}(b\varpi_{n-1}) \rightarrow \mathscr{D}(\mathscr{H}_{n,m})_{q,\mathfrak{k}},
\\
t_1^{i_1}\cdots t_n^{i_n} \boxtimes t_{n+1}^{*i_{n+1}}\cdots t_N^{*i_N} \otimes t_N^{j_N}\cdots
t_{n+1}^{j_{n+1}} \boxtimes t_n^{*j_n}\cdots t_1^{*j_1}\\
\qquad{} \mapsto t_1^{i_1}\cdots t_n^{i_n}
t_{n+1}^{*i_{n+1}}\cdots t_N^{*i_N}f_0t_N^{j_N}\cdots t_{n+1}^{j_{n+1}}t_n^{*j_n}\cdots t_1^{*j_1}.
\end{gather*}
Recall that, as in the classical case, see \cite{VO},
\[
L^{(n)}(a\varpi_1) \otimes L^{(n)}(b\varpi_{n-1}) = \oplus_{i=0}^{\max(a,b)}
L^{(n)}((a-i)\varpi_1+(b-i)\varpi_{n-1})
\]
in the category of $U_q \mathfrak{sl}_n$-modules. By explicit calculations one can show that the $U_q
\mathfrak{sl}_n$-highest weight vector in $L^{(n)}((a-i)\varpi_1+(b-i)\varpi_{n-1}) \subset
L^{(n)}(a\varpi_1) \otimes L^{(n)}(b\varpi_{n-1})$ has the form
\[
t_1^{a-i}  \left(\sum_{j=1}^nt_j
\otimes t_j^*\right)^{i}   t_n^{*(b-i)}.
\]

Now a routine application of the commutation relations (similar to the ones in Remark~\ref{remark2}) allows one to reduce
every highest weight vector to a linear span of the vectors $t_1^at_N^{*b}\varphi(x_{n+1})t_{n+1}^ct_n^{*d}$,
where $\varphi(x_{n+1})$ is a f\/inite function.
\end{proof}



Now let us obtain an explicit form of the restriction of $\nu_q$ to the space
$\mathscr{D}(\mathscr{H}_{n,m})_{q,\mathfrak{k}}^{U_q \mathfrak{k}}$ of $U_q \mathfrak{k}$-invariant elements
of $\mathscr{D}(\mathscr{H}_{n,m})_{q,\mathfrak{k}}$.

Recall the standard notation $(a;q)_k=(1-a)\cdots (1-aq^{k-1})$, $(a;q)_0=1$,
\[
\int_{1}^\infty f(x) d_{q^{-2}}x=\big(q^{-2}-1\big)\sum_{k=0}^\infty f\big(q^{-2k}\big)q^{-2k}.
\]

\begin{proposition}\label{rad_int}
For any function $f(x_{n+1}) \in \mathscr{D}(\mathscr{H}_{n,m})_{q,\mathfrak{k}}$ one has
\[
\int_{\mathscr{H}_{n,m}} f(x_{n+1}) d\nu_q= \int_1^\infty f(x)\rho(x) d_{q^{-2}}x,
\]
where
\begin{gather}\label{rho_x}
\rho(x)=\mathrm{const_2}\,   x^{m-1}\big(q^{-2}x-1\big)\big(q^{-4}x-1\big)\cdots \big(q^{-2(n-1)}x-1\big),
\end{gather}
and $\mathrm{const_2}=\frac{1}{q^{-2}-1}\prod \limits_{j=1}^{n-1}\frac{1}{q^{-2j}-1}.$
\end{proposition}

\begin{proof} By explicit calculations,
\begin{gather*}
\int_{\mathscr{H}_{n,m}} f(x_{n+1}) d\nu_q= \mathrm{const_1}
\sum_{\substack{i_1\ldots,i_n \in -\mathbb{Z}_+,\\
i_{n+1},\ldots,i_{N-1} \in \mathbb{N}}} f\big(q^{2i_1+\cdots+2i_n}\big)
q^{2(N-1)i_1+2(N-2)i_2+\cdots+2i_{N-1}}
\\
\qquad{} = \mathrm{const_1}    \sum \limits_{i_{n+1},\ldots,i_{N-1}=1}^\infty \left(\sum_{\substack{i_1, \ldots, i_n
\in -\mathbb Z_+, \\ i_1+\cdots+i_n=i}} f(q^{2i})
q^{2(N-1)i_1+\cdots+2mi_n}\right) q^{2(m-1)i_{n+1}+\cdots+2i_{N-1}}
\\
\qquad{}= \mathrm{const_1} \sum_{i \in -\mathbb Z_+} \frac{q^{m(m-1)}}{(q^2;q^2)_{m-1}} \left(\sum_{\substack{i_1, \ldots, i_{n-1}
\in -\mathbb Z_+, \\ i_1+\cdots+i_{n-1} \geq i}}
q^{2(n-1)i_1+\cdots+2i_{n-1}}\right)f(q^{2i})q^{2mi}
\\
\qquad{}= \mathrm{const_1} \frac{q^{m(m-1)}}{(q^2;q^2)_{m-1}} \sum_{i \in -\mathbb Z_+} f(q^{2i})q^{2mi}  \left(
\sum_{\substack{i_1, \ldots, i_{n-1} \in \mathbb Z_+, \\ i_1+\cdots+i_{n-1} \leq -i}}
q^{-2(n-1)i_1-\cdots-2i_{n-1}}\right).
\end{gather*}
Let us verify that for every $j,k \in \mathbb Z_+$
\[
\sum_{\substack{i_1, \ldots, i_k \in \mathbb Z_+, \\ i_1+\cdots+i_k
\leq j}}
q^{-2(k-1)i_1-\cdots-2i_k}=\frac{(q^{-2};q^{-2})_{j+k}}{(q^{-2};q^{-2})_j(q^{-2};q^{-2})_k}.
\]
Denote the l.h.s.\ of the previous equation by $\Psi(j,k)$. One can verify the recurrence relation for the
$q$-Pascal triangle
\[
\Psi(j,k)=q^{-2k}\Psi(j-1,k)+\Psi(j,k-1)
\]
and the boundary values
\[
\Psi(0,k)=1, \qquad \Psi(j,1)=\frac{1-q^{-2(j+1)}}{1-q^{-2}},
\]
explicitly. Thus $\Psi(j,k)$ are the corresponding entries of the $q$-Pascal triangle. Now we can complete our
calculations and obtain that
\begin{gather*}
\int_{\mathscr{H}_{n,m}} f(x_{n+1}) d\nu_q= \mathrm{const_1}
\frac{q^{m(m-1)}}{(q^2;q^2)_{m-1}} \sum_{i \in -\mathbb Z_+} f\big(q^{2i}\big)
q^{2mi}
\frac{(q^{-2};q^{-2})_{n-1-i}}{(q^{-2};q^{-2})_{-i}(q^{-2};q^{-2})_{n-1}}
\\
\qquad {} = \mathrm{const_1}   \frac{q^{m(m-1)}}{(q^2;q^2)_{m-1}} \sum_{i
\in -\mathbb Z_+} f\big(q^{2i}\big) q^{2mi}
\frac{(q^{2i-2};q^{-2})_{n-1}}{(q^{-2};q^{-2})_{n-1}}=\int_1^\infty
f(x)\rho(x) d_{q^{-2}}x,
\end{gather*}
where $\rho(x)=\mathrm{const_2} \,  x^{m-1}(q^{-2}x-1)(q^{-4}x-1)\cdots (q^{-2(n-1)}x-1),$ and
\begin{gather*}
\mathrm{const_2}=\frac{1}{q^{-2}-1} \mathrm{const_1}   \frac{q^{m(m-1)}}{(q^2;q^2)_{m-1}}
\frac{(-1)^{n-1}}{(q^{-2};q^{-2})_{n-1}} =
\frac{(-1)^{n-1}}{(q^{-2}-1)(q^{-2};q^{-2})_{n-1}}.  \tag*{\qed}
\end{gather*}
\renewcommand{\qed}{}
\end{proof}

Now one can complete $\mathscr{D}(\mathscr{H}_{n,m})_{q,\mathfrak{k}}^{U_q \mathfrak{k}}$ with respect to the
norm $||f||^2=\int_1^\infty f^*f \rho(x)d_{q^{-2}}x$. The resulting Hilbert space will be denoted by
$L^2\big(d\nu_q^{(0)}\big)$.



\section[Covariant first order differential calculi over $\operatorname{Pol}(\widetilde{\mathscr{H}}_{n,m})_q$]{Covariant f\/irst order dif\/ferential calculi over $\boldsymbol{\operatorname{Pol}(\widetilde{\mathscr{H}}_{n,m})_q}$}
\label{sec_5}

In this section we introduce   holomorphic and antiholomorphic covariant f\/irst order dif\/ferential calculi
over $\operatorname{Pol}(\widetilde{\mathscr{H}}_{n,m})_q$.

First of all, recall a general def\/inition of a covariant f\/irst order dif\/ferential calculus fol\-lowing~\cite{Kl-Sch}.

Let $F$ be a unital algebra. A f\/irst order dif\/ferential calculus over $F$ is a pair $(M,d)$ where $M$ is an
$F$-bimodule and $d:F \rightarrow M$ is a linear map such that
\begin{enumerate}\itemsep=0pt
\item for all $f_1,f_2 \in F$ one has $d(f_1 \cdot f_2)=df_1 \cdot f_2+f_1 \cdot df_2$;
\item $M$ is a linear span of the vectors $f_1 \cdot df_2 \cdot f_3$, where $f_1,f_2,f_3 \in F$.
\end{enumerate}

Now suppose that $A$ is a Hopf algebra and $F$ is an $A$-module algebra. A f\/irst order dif\/ferential calculus
$(M,d)$ over $F$ is called covariant if the following conditions hold:
\begin{enumerate}\itemsep=0pt
\item $M$ is an $A$-covariant $F$-bimodule, i.e.\ the action maps $F \otimes M \rightarrow M$, $M \otimes F \rightarrow M$
are morphisms of $A$-modules;
\item $d$ is a morphism of $A$-modules.
\end{enumerate}

Let $\mathcal{L}=L \circ \omega$, where $L$ is the canonical antirepresentation of $U_q \mathfrak{sl}_N$ in
$\mathbb C[SL_N]_q$ (or the representation of $U_q \mathfrak{sl}_N^{\rm op}$) and $\omega$ is an
antiautomorphism of $U_q \mathfrak{sl}_N$ def\/ined on the generators by $\omega(K_i)=K_i^{-1}$,
$\omega(E_i)=F_i$, $\omega(F_i)=E_i$. Thus $\mathcal{L}$ is a representation of $U_q \mathfrak{sl}_N$ in
$\mathbb C[SL_N]_q$.

Consider the action of $\mathcal{L}(E_1K_1^{-1/2})$ in $\mathbb C [SL_N]_q$ and its restriction to
$\operatorname{Pol}(\widetilde{\mathscr{H}}_{n,m})_q$. Put
\[\partial: \
\operatorname{Pol}(\widetilde{\mathscr{H}}_{n,m})_q \rightarrow \operatorname{Pol}(\widetilde{X})_q, \qquad \partial
t_i=\mathcal{L}\big(E_1K_1^{-1/2}\big)t_{1i}.
\]

Let $\Omega^{(1,0)}(\widetilde{\mathscr{H}}_{n,m})_q \subset \operatorname{Pol}(\widetilde{X})_q$ be the
$\operatorname{Pol}(\widetilde{\mathscr{H}}_{n,m})_q$-submodule generated by $\partial t_i$, $i=1,\dots,N$.
\begin{lemma}
$\Omega^{(1,0)}(\widetilde{\mathscr{H}}_{n,m})_q$ is a $U_q \mathfrak{sl}_N$-covariant first order differential calculus
over $\operatorname{Pol}(\widetilde{\mathscr{H}}_{n,m})_q$.
\end{lemma}
\begin{proof} One has to verify the Leibniz rule  which immediately follows from the formulas for the
comultiplication in $U_q \mathfrak{sl}_N$. Since left and right action of $U_q \mathfrak{sl}_N$ in $\mathbb C
[SL_N]_q$ commute, $\partial$ is a~morphism of $U_q \mathfrak{sl}_N$-modules, so the calculus is covariant.
\end{proof}

Now we def\/ine $\bar{\partial}: \operatorname{Pol}(\widetilde{\mathscr{H}}_{n,m})_q \rightarrow
\operatorname{Pol}(\widetilde{X})_q$ by the rule $\bar{\partial}f=(\partial f^*)^*$. Let
$\Omega^{(0,1)}(\widetilde{\mathscr{H}}_{n,m})_q \subset \operatorname{Pol}(\widetilde{X})_q$ be the
$\operatorname{Pol}(\widetilde{\mathscr{H}}_{n,m})_q$-submodule generated by $\bar{\partial} t_i$, $i=1,\dots,N$.
\begin{lemma}
$\Omega^{(0,1)}(\widetilde{\mathscr{H}}_{n,m})_q$ is a $U_q \mathfrak{sl}_N$-covariant first order differential calculus
over $\operatorname{Pol}(\widetilde{\mathscr{H}}_{n,m})_q$.
\end{lemma}

This lemma can be proved similarly to the previous one.
\begin{remark}
The introduced f\/irst order dif\/ferential calculi can be obtained in another
way. One should start from the canonical Wess--Zumino calculi on a quantum
complex space introduced in~\cite{RTF} and then turn to a localization of
the corresponding algebras of functions.
\end{remark}



Now we introduce a Hermitian pairing $\Omega^{(0,1)}(\widetilde{\mathscr{H}}_{n,m})_q \times
\Omega^{(0,1)}(\widetilde{\mathscr{H}}_{n,m})_q \rightarrow \operatorname{Pol}({\mathscr{H}}_{n,m})_q$.

Let $P: \operatorname{Pol}(\widetilde{X})_q \rightarrow \operatorname{Pol}({\mathscr{H}}_{n,m})_q$ be the
projection parallel to a sum of other $U_q \mathfrak{s}(\mathfrak{gl}_1 \times \mathfrak{gl}_{N-1})$-isotypic components of $\mathcal{L}$. Now we def\/ine
\[
(\theta_1,\theta_2)=P(\theta^*_2\theta_1), \qquad \theta_1,\theta_2 \in \Omega^{(0,1)}(\widetilde{\mathscr{H}}_{n,m})_q.
\]
By obvious commutativity of the left and right actions of $U_q \mathfrak{sl}_N$ we have

\begin{proposition}
The pairing $(\cdot,\cdot)$ is $U_q \mathfrak{sl}_N$-invariant.
\end{proposition}
\begin{lemma}
\[P(t_{2j}t^*_{2k})=q^{-2}\frac{1-q^2}{1-q^{2(N-1)}}(\varepsilon_{jk}-t_{1j}t^*_{1k}),\] where
\[
\varepsilon_{jk}=\begin{cases} q^{2(j-1)}, \quad & j=k, \  j \geq n+1, \\
-q^{2(j-1)}, \quad & j=k, \  j \leq n, \\
0, \quad & j \neq k.
\end{cases}
\]
\end{lemma}
\begin{proof} Since $U_q \mathfrak{s(gl}_1 \times \mathfrak{gl}_{N-1})=\mathbb C[K_1,K_1^{-1}] \otimes U_q
\mathfrak{sl}_{N-1}$, one can decompose the projection $P$ as follows: $P=P_0P_1$, where $P_1$ is a
projection to the subspace of $U_q \mathfrak{sl}_{N-1}$-invariant elements (w.r.t.\ the $\mathcal{L}$-action)
and $P_0$ is a projection to the subspace of elements that are preserved by the $\mathcal{L}(K_1)$-action.

Let $u_1,\dots,u_k$ be the standard basis in the $U_q \mathfrak{sl}_k$-module $L(\varpi_1)$, and
$v_1,\ldots,v_k$ the dual basis in the $U_q \mathfrak{sl}_k$-module $L(\varpi_{k-1})$, where $\varpi_1$ and
$\varpi_{k-1}$ are the fundamental weights. A standard argument on f\/inite dimensional $U_q
\mathfrak{sl}_k$-modules allows one to prove that $\sum\limits_{j=1}^k (-q)^{j-1}u_j \otimes v_j$ is a~$U_q
\mathfrak{sl}_k$-invariant element in the $U_q \mathfrak{sl}_k$-module $L(\varpi_1) \otimes L(\varpi_{k-1})$,
and the map
\begin{gather}\label{proj_inv}
u_i \otimes v_j \mapsto \begin{cases} \displaystyle \frac{1-q^2}{1-q^{2k}}(-q)^{i-1}\sum_{a=1}^k (-q)^{a-1} u_a \otimes v_a, \quad &  i=j,
\\ 0, \quad & i \neq j,
\end{cases}
\end{gather}
is a projection to the subspace of $U_q \mathfrak{sl}_k$-invariant elements parallel to other isotypic
components. By obvious reasons, for $j=1,\ldots,N-1$ the maps
\[
\phi_j: \ u_i \mapsto t_{i+1,j},\qquad \psi_j: \ v_i \mapsto \det \nolimits_q T_{i+1,j},
 \] admit extensions to morphisms of
$U_q \mathfrak{sl}_{N-1}$-modules
\[
\phi_j: \ L(\varpi_{N-2}) \rightarrow \mathbb C[SL_N]_q, \qquad \psi_j: \ L(\varpi_1) \rightarrow \mathbb
C[SL_N]_q
\]
(w.r.t.\ the $\mathcal{L}$-action). By def\/inition~\eqref{def_inv}, $t_{ij}^*=(-q)^{j-i}\det \nolimits_qT_{ij}$
for $i \leq m$ and $j \geq n+1$. Thus one can apply the map \eqref{proj_inv} to compute for $l \geq n+1$
\[
P_1(t_{2l}t_{2l}^*)=(-q)^{l-2}P_1(t_{2l} \det \nolimits_q T_{2l})
=(-q)^{l-2}\frac{1-q^2}{1-q^{2(N-1)}}\sum_{j=1}^{N-1}(-q)^{j-1}t_{j+1,l}\det \nolimits_q T_{j+1,l}.
\]
Since $\mathbb C[SL_N]_q$ is a Hopf algebra, one has $\sum\limits_{j=1}^{N} t_{lj}S(t_jk)=\delta_{lk}$ and
$S(t_{jk})=(-q)^{j-k}\det \nolimits_q T_{kj}$. Thus
\begin{gather*}
P_1(t_{2l}t_{2l}^*)=(-q)^{2(l-2)}\frac{1-q^2}{1-q^{2(N-1)}}\big( 1-(-q)^{1-l}t_{1l}\det \nolimits_q
T_{1l}\big)
\\
\phantom{P_1(t_{2l}t_{2l}^*)}{} = (-q)^{2(l-2)}\frac{1-q^2}{1-q^{2(N-1)}}\big( 1-(-q)^{-2(l-1)}t_{1l}t_{1l}^*\big).
\end{gather*}
The other cases can be verif\/ied in a similar way. \end{proof}

\begin{proposition}
\[
\big(\bar \partial x_{n+1}, \bar \partial x_{n+1}\big)= q^{2(n-1)} \frac{1-q^2}{1-q^{2(N-1)}}
x_{n+1}(1-q^{-2n}x_{n+1}).
\]
\end{proposition}

\begin{proof} An easy application of the previous lemma allows one to compute that
\[
\big(\bar \partial (t_{1k}^*t_{1l}), \bar \partial
(t_{1j}^*t_{1i})\big)=q^{-2}\frac{1-q^2}{1-q^{2(N-1)}}(\varepsilon_{jk}t_{1i}^*t_{1l}-(t_{1i}^*t_{1j})(t_{1k}^*t_{1l})).
\]
Thus for $j,k \geq n+1$ one has
\begin{gather*}
\big(\bar \partial (q^{-2j}t_{1j}^*t_{1j}), \bar \partial
(q^{-2k}t_{1k}^*t_{1k})\big)\\
\qquad{} =q^{-2}\frac{1-q^2}{1-q^{2(N-1)}}\big(\delta_{jk}q^{-2(j+1)}t_{1j}^*t_{1j}
-\big(q^{-2j}t_{1j}^*t_{1j}\big)\big(q^{-2k}t_{1k}^*t_{1k}\big)\big).
\end{gather*}
By easy computations, $q^{-2(n+1)}x_{n+1}=\sum\limits_{j=n+1}^N q^{-2j}t_{1j}^*t_{1j}$, which enables to prove the
claim.
\end{proof}

Let us f\/ix notation for $q$-dif\/ference operators:
\[
B_-: \ f(x) \mapsto \frac{f(q^{-2}x)-f(x)}{q^{-2}x-x}, \qquad B_+: \ f(x) \mapsto \frac{f(q^2x)-f(x)}{q^2x-x}.
\]
\begin{lemma}\label{lem_dbar}
\[
\big(\bar \partial f(x_{n+1}), \bar \partial g(x_{n+1})\big)= q^{2(n-1)} \frac{1-q^2}{1-q^{2(N-1)}} x_{n+1}
\big(1-q^{-2n}x_{n+1}\big) \overline{B_-g(x_{n+1})} B_-f(x_{n+1}).
\]
\end{lemma}
\begin{proof} By explicit calculations in $\mathbb C[SL_N]_q$ we obtain for $i \geq n+1$
\[
t_{1i}\bar \partial x_{n+1}= q^{-1} \bar \partial x_{n+1} t_{1i}, \qquad t_{1i}^* \bar \partial
x_{n+1}=q^{-1} \bar
\partial x_{n+1} t_{1i}^*.
\]
Let us verify the second identity. One has $\bar \partial x_{n+1}=\sum\limits_{j=n+1}^Nt_{1j}t_{2j}^*$, so
\begin{gather*}
t_{1i}^*\bar \partial x_{n+1}=t_{1i}^*  \left(\sum_{j=n+1}^{i-1}
t_{1j}t_{2j}^*+t_{1i}t_{2i}^*+\sum_{j=i+1}^Nt_{1j}t_{2j}^*\right) = q^{-1} \sum_{j=n+1}^{i-1}t_{1j}t_{2j}^*
  t_{1i}^*
\\
{}
+q^{-1}t_{1i}t_{2i}^*t_{1i}^*+\big(q-q^{-1}\big)\sum_{j>i}t_{1j}t_{1j}^*t_{2i}^*+q^{-1}\sum_{j=i+1}^N
t_{1j}t_{2j}^*t_{1i}^*+\big(q^{-1}-q\big)\sum_{j>i}t_{1j}t_{1j}^*t_{2i}^*.
\end{gather*}
So we have $x_{n+1}\bar \partial x_{n+1}=q^{-2} \bar
\partial x_{n+1} x_{n+1}$. Hence for every polynomial $f(x)$
\[
\bar \partial f(x_{n+1})=\bar \partial x_{n+1}\frac{f(q^{-2}x_{n+1})-f(x_{n+1})}{q^{-2}x_{n+1}-x_{n+1}}=\bar
\partial x_{n+1} B_-(f)(x_{n+1}),
\]
and  the claim follows from the previous proposition. \end{proof}

Using the above formal arguments, we
can extend $\bar \partial$ to
$\mathscr{D}(\mathscr{H}_{n,m})_{q,\mathfrak{k}}$.


\section[The Laplace-Beltrami operator and its radial part]{The Laplace--Beltrami operator and its radial part}\label{sec_6}

In this section we introduce a $U_q \mathfrak{sl}_N$-invariant operator $\square$ in the space of f\/inite
functions $\mathscr{D}(\mathscr{H}_{n,m})_{q, \mathfrak{k}}$. This operator will be considered as a quantum
analog for the invariant Laplace--Beltrami operator on the complex hyperbolic space $\mathscr{H}_{n,m}$. Also
we compute an explicit formula for the restriction $\square^{(0)}$ of $\square$ to the space
$\mathscr{D}(\mathscr{H}_{n,m})_{q,\mathfrak{k}}^{U_q \mathfrak{k}}$, the so called radial part of~$\square$.
The restriction appears to be a $q$-dif\/ference operator in variable $x=x_{n+1}$.

Now we def\/ine $\square$ by the formula
\[
\int_{\mathscr{H}_{n,m}}f_2^* (\square f_1) d\nu_q=\int _{\mathscr{H}_{n,m}} (\bar \partial f_1, \bar
\partial f_2) d\nu_q, \qquad f_1,f_2 \in \mathscr{D}(\mathscr{H}_{n,m})_{q,\mathfrak{k}}.
\]

\begin{proposition}
$\square$ is a self-adjoint $U_q \mathfrak{sl}_N$-invariant operator.
\end{proposition}

\begin{proof} The self-adjointness follows from the def\/inition. The $U_q \mathfrak{sl}_N$-invariance of the
form $(\cdot,\cdot)$ and the linear functional $\int_{\mathscr{H}_{n,m}} \cdot d\nu_q$ implies the $U_q
\mathfrak{sl}_N$-invariance of~$\square$.
\end{proof}



Let us f\/ind the restriction $\square^{(0)}$ of the operator $\square$ on the subspace
\[
\mathscr{D}(\mathscr{H}_{n,m})_{q,\mathfrak{k}}^{U_q \mathfrak{k}}=\big\{ f(x_{n+1}), \ \mathrm{supp} f \subset q^{-2 \mathbb{Z}_+},\
\sharp (\mathrm{supp} f) < \infty \big\},
\]
based on the equation
\begin{gather*}
\int_1^\infty \big(\square^{(0)} f\big) (x) \overline{f(x)} \rho(x) d_{q^{-2}} x \\
\qquad{} = q^{2(n-1)}
\frac{1-q^2}{1-q^{2(N-1)}} \int_1^\infty x \big(1-q^{-2n} x\big) \left| \frac{f(q^{-2}x)-f(x)}{q^{-2}x - x}
\right|^2 \rho(x) d_{q^{-2}} x
\end{gather*}
for every $f$ with f\/inite support.

Using the $q$-analog of the partial integration formulas one can prove that operators $B_-$ and $-q^2 B_+$
are formally dual. Exactly, for every functions $u(x)$, $v(x)$ with f\/inite support on $q^{2\mathbb{Z}}$ holds
\[ \int_0^\infty u(x) \frac{v(q^{-2}x)-v(x)}{(q^{-2}-1)x} d_{q^{-2}} x = -q^2 \int_0^\infty
\frac{u(x)-u(q^2x)}{(1-q^2)x} v(x) d_{q^{-2}} x,
\]
where
\[
\int_0^\infty f(x) d_{q^{-2}} x = \big(q^{-2} - 1\big)
\sum\limits_{k=-\infty}^\infty f\big(q^{-2k}\big) q^{-2k}.
\]

Thus
\[
\square^{(0)}: \ f(x) \mapsto \mathrm{const}(q, n, N) \rho(x)^{-1} B_+ x \big(q^{-2n} x - 1\big) \rho(x) B_- f(x),
\]
where
\[
\mathrm{const}(q, n, N) = q^{2n} \frac{1-q^2}{1-q^{2(N-1)}},
\]
and $\rho(x)$ is def\/ined in \eqref{rho_x}.

Lemma~\ref{lem_dbar} allows us to extend the Hermitian $U_q
\mathfrak{sl}_N$-invariant pairing to f\/irst order dif\/ferential forms with
coef\/f\/icients in $\mathscr{D}(\mathscr{H}_{n,m})_{q,\mathfrak{k}}$.

\begin{lemma}
For every function $f(x)$ with finite support on $q^{-2\mathbb{Z}_+}$
\begin{gather*}
\square^{(0)} f(x) = \frac{q^2}{(1-q^2)^2x} \big((x-1)q^{2m-2}f(q^2x) +
\big(q^{-2n}x-1\big)f\big(q^{-2}x\big)\\
\phantom{\square^{(0)} f(x) =}{}
+\big(1+q^{2m-2}-q^{2m-2}x-q^{-2n}x\big)f(x)\big).
\end{gather*}
\end{lemma}

\begin{proposition}\label{cont_spec}
The operator $\square^{(0)}$ is bounded.
\end{proposition}

\begin{proof} Consider the operator
\[
\widetilde{\square} : \ f(x) \mapsto x^{-(N-1)} \mathrm{const}(q, n, N) q^{-2n} B_+ x^{N+1} B_-f(x).
\]
It dif\/fers from $\square^{(0)}$ by a compact operator, so it is suf\/f\/icient to prove that
$\widetilde{\square}$ is bounded. The boundness of the latter operator can be proved by the direct
evaluation.
\end{proof}

By the previous proposition, one can extend $\square^{(0)}$ from
$\mathscr{D}(\mathscr{H}_{n,m})_{q,\mathfrak{k}}^{U_q \mathfrak{k}}$ to a bounded self-adjoint operator in
$L^2\big(d\nu_q^{(0)}\big)$. The latter extension will also be denoted by $\square^{(0)}$.


\section[Generalized eigenfunctions of $\square^{(0)}$ and Al-Salam-Chihara polynomials]{Generalized eigenfunctions of $\boldsymbol{\square^{(0)}}$\\ and Al-Salam--Chihara polynomials}\label{sec_7}

In this section we obtain the initial results on the bounded self-adjoint operator $\square^{(0)}$, namely we
obtain its formal eigenfunctions and eigenvalues explicitly. Note that explicit computations of the
asymptotics of these eigenfunctions, as in the classical case, allow us to consider a quantum analog of the
Harish-Chandra $c$-function (see Appendix~\ref{appendixA}).

By the direct computation we obtain the following lemma.

\begin{lemma}\label{c1.5.2} The function
\[
\Phi_l(x) = {_{3}\Phi_{2}}\left(
\begin{array}{c}
x,q^{-2l},q^{2(l+N-1)}\\ q^{2n},0
\end{array}
;q^2,q^2\right)
\]
in $\mathscr{D}(\mathscr{H}_{n,m})_q^\prime$ is a generalized eigenfunction for $\square^{(0)}$:
\[\square^{(0)} \Phi_l=\lambda(l)\Phi_l\] with the eigenvalue
\[
\lambda(l) = - q^{2-2n}\frac{(1-q^{-2l})(1-q^{2l+2(N-1)})}{(1-q^2)^2}.
\]
\end{lemma}



Recall the def\/inition of the Al-Salam--Chihara polynomials, following \cite{koekoek}:
\[
Q_k(z;a,b|q)=\frac{(ab;q)_k}{a^k}\ {_{3}\Phi_{2}}\left(
\begin{array}{c}
q^{-k},a e^{i\theta},ae^{-i\theta}\\ ab,0
\end{array}
;q,q\right), \qquad z=\cos\theta.
\]

We have the following well known orthogonality relations for the Al-Salam--Chihara polynomials:
\begin{gather}
\frac{1}{2\pi}\int_{-1}^1 Q_i(z;a,b|q)Q_j(z;a,b|q)\frac{w(z)}{\sqrt{1-z^2}}\,dz\nonumber\\
\qquad{}+\sum_{1<aq^k<a}w_k Q_i(z_k;a,b|q)Q_j(z_k;a,b|q)=\frac{\delta_{ij}}{(q^{i+1},abq^i;q)_\infty},\label{rel_orth_ASC}
\end{gather}
where
\[
w(z)=\frac{h(z,1)h(z,-1)h(z,q^{1/2})h(z,-q^{1/2})}{h(z,a)h(z,b)}, \!\!\qquad
h(z,a)=\big(ae^{i\theta},ae^{-i\theta};q\big)_{\infty},\!\! \qquad z=\cos\theta,
\]
$z_k=\frac{aq^k+aq^{-k}}{2}$, and
\[
w_k=\frac{(a^{-2};q)_\infty}{(q,ab,b/a;q)_\infty}\frac{(1-a^2q^{2k})(a^2,ab;q)_k}
{(1-a^2)(q,qa/b;q)_k}q^{-k^2}\left(\frac{1}{a^3b}\right)^k.
\]

Also, these polynomials satisfy the following recurrence relations:
\begin{gather}
zQ_i(z;a,b|q)=\frac12
Q_{i+1}(z;a,b|q)+\frac12q^i(a+b)Q_i(z;a,b|q)\nonumber\\
\phantom{zQ_i(z;a,b|q)=}{} +\frac12\big(1-q^i\big)\big(1-abq^{i-1}\big)Q_{i-1}(z;a,b|q).\label{rec_ASC}
\end{gather}

As one can see, the eigenfunctions $\Phi_l(x)$ are connected with the Al-Salam--Chihara polynomials:
\[
\Phi_l(q^{-2k}) = \frac{q^{k(N-1)}}{(q^{2n};q^2)_k}Q_k\big(z;q^{2n-N+1},q^{N-1}|q^2\big), \qquad
e^{i\theta}=q^{2l+N-1},\qquad z=\cos\theta.
\]

Let us denote the orthogonal measure for Al-Salam--Chihara polynomials by $d\sigma$.  Note that if $2n-N+1>0$,
the sum in the orthogonality relations vanishes and we have just the continuous measure.



\section[A spectral theorem for the radial part of $\square$]{A spectral theorem for the radial part of $\boldsymbol{\square}$}\label{sec_9}

In this section we obtain a spectral theorem for $\square^{(0)}$. As in the classical case~\cite[pp.~429--432]{Faraut}, the support of the Plancherel measure consists of continuous and discrete parts.
The continuous part corresponds to principal unitary series of $U_q \mathfrak{su}_{n,m}$-modules related to a
quantum analog of the cone
\[
\Xi_{n,m}=\big\{x \in \mathbb C^{N}\,|\, {-}x_1 \bar{x}_1-\cdots -x_n \bar{x}_n + x_{n+1}\bar{x}_{n+1}+\cdots
+x_N\bar{x}_N=0\big\}
\]
(this series is established in~\cite{BS}). The discrete part is supposed to be corresponded to a discrete
series of unitary $U_q \mathfrak{su}_{n,m}$-modules.

\begin{theorem}
The bounded self-adjoint linear operator $\square^{(0)}$ is unitary equivalent to the operator of
multiplication by independent variable in the Hilbert space $L^2(d\sigma)$. The unitary equivalence is given
by the operator
\begin{gather*}
U:\  L^2\big(d\nu_q^{(0)}\big)  \rightarrow L^2(d\sigma),
\\
U:\ f(x) \mapsto\hat{f}(\lambda)=\int_1^\infty f(x)\Phi_l(x)\rho(x)d_{q^{-2}}x,
\end{gather*}
where $\lambda=-q^{2-2n}\frac{(1-q^{-2l})(1-q^{2l+2N-2})}{(1-q^2)^2}$.
\end{theorem}

\begin{proof} Let us consider the f\/inite functions on $q^{-2\mathbb Z_+}$
\[
f_j(x)=
\begin{cases} 1, & x=q^{-2j},\\ 0, & \text{otherwise}, \end{cases}
\qquad j \in \mathbb Z_+.
\]
These functions form an orthogonal system in $\mathscr{D}(\mathscr{H}_{n,m})_{q,\mathfrak{k}}^{U_q \mathfrak{k}}$
and the completion of its linear span is $L^2\big(d\nu_q^{(0)}\big)$. By standard arguments~\cite{AG}, the bounded
self-adjoint linear operator $\square^{(0)}$ is unitary equivalent to the multiplication operator $f(\lambda)
\mapsto \lambda f(\lambda)$ in the Hilbert space $L^2(d\mu(\lambda))$ of square integrable functions with
respect to a certain measure $d\mu(\lambda)$ with compact support in $\mathbb R$. Let us f\/ind explicitly the
corresponding measure and the operator of unitary equivalence $U$. One can f\/ix the unitary equivalence
operator by the condition $Uf_0=1$.

By easy calculations, $\square^{(0)}$ transforms $f_j$ by the formula
\begin{gather*}
\square^{(0)}f_j=\frac{q^2}{(1-q^2)^2} \big( f_{j+1}\big(1-q^{2j+2}\big)q^{2m-2}+f_{j-1}\big(1-q^{2j+2n-2}\big)q^{-2n}\\
\phantom{\square^{(0)}f_j=}{} +f_j
\big(q^{2j}+q^{2j+2m-2}-q^{2m-2}-q^{-2n}\big)\big), \qquad j \in \mathbb Z_+,
\end{gather*}
and $||f_j||^2=q^{-2j(N-1)}\frac{(q^{2j+2};q^2)_{n-1}}{(q^2;q^2)_{n-1}}$ (here we naturally suppose
$f_{-1}=0$).

{\sloppy Thus the f\/inite functions $e_j=q^{j(N-1)}\sqrt{\frac{(q^{2};q^2)_{n-1}}{(q^{2j+2};q^2)_{n-1}}}f_j$ form an
orthonormal system in $L^2\big(d\nu_q^{(0)}\big)$ and $\square^{(0)}$ acts on them by the formula
\begin{gather*}
\square^{(0)}e_j=\frac{q^{m-n+1}}{(1-q^2)^2}\Big(
e_{j+1}\sqrt{\big(1-q^{2j+2}\big)\big(1-q^{2j+2n}\big)}+e_{j-1}\sqrt{\big(1-q^{2j}\big)\big(1-q^{2j+2n-2}\big)} \\
\phantom{\square^{(0)}e_j=}{}+e_j
\big(q^{2j+(N-1)}+q^{2j+2n-(N-1)}-q^{N-1}-q^{1-N}\big)\Big), \qquad j \in \mathbb Z_+.
\end{gather*}

}

Thus $P_j=Ue_j \in L^2(d\mu(\lambda))$, $j\in\mathbb{Z}_+$ form an orthonormal system of polynomials
\[
\int P_i(\lambda)P_j(\lambda)d\mu(\lambda)=\delta_{ij},\qquad i,j\in\mathbb{Z}_+,
\]
and one has{\samepage
\begin{gather}
\lambda P_j(\lambda)= \frac{q^{m-n+1}}{(1-q^2)^2}
\Big(P_{j+1}(\lambda)\sqrt{\big(1-q^{2j+2}\big)\big(1-q^{2j+2n}\big)}\nonumber\\
\phantom{\lambda P_j(\lambda)=}{}
 +P_{j-1} (\lambda) \sqrt{\big(1-q^{2j}\big)
\big(1-q^{2j+2n-2}\big)}\nonumber\\
\phantom{\lambda P_j(\lambda)=}{}
+P_j(\lambda) \big(q^{2j+(N-1)}+q^{2j+2n-(N-1)}-q^{N-1}-q^{1-N}\big)\Big), \qquad j \in \mathbb Z_+,\label{P_rec}
\\
\label{init_P0}
P_0(\lambda)=Ue_0=Uf_0=1.
\end{gather}
(we naturally suppose that $P_{-1}=0$).}

The orthogonal polynomials $P_j(\lambda)$, $j\in\mathbb{Z}_+$ are determined by~\eqref{P_rec} and
\eqref{init_P0}. Let us compare them with the corresponding recurrence relations~\eqref{rec_ASC} and the
initial data for the Al-Salam--Chihara polynomials $Q_j(z;q^{2n-(N-1)},q^{N-1}|q^2)$. One can observe that
\[
||Q_j||^2=\frac{1}{(q^{2j+2},q^{2j+2n};q^2)_\infty}=\frac{(q^2,q^{2n};q^2)_j}{(q^2;q^2)_{n-1}(q^{2n};q^2)^2_\infty}.
\]

The polynomials
\[
\widetilde{Q}_j\overset{\mathrm{def}}{=}\sqrt{\frac{1}{(q^2,q^{2n};q^2)_j}} \,Q_j
\]
satisfy the following recurrence relations:
\begin{gather*}
z\widetilde{Q}_j=\dfrac12\sqrt{\big(1-q^{2(j+1)}\big)\big(1-q^{2j+2n}\big)}\widetilde{Q}_{j+1}+
\dfrac12\sqrt{\big(1-q^{2j}\big)\big(1-q^{2j+2n-2}\big)}\widetilde{Q}_{j-1}\\
\phantom{z\widetilde{Q}_j=}{} +\dfrac12q^{2j}\big(q^{2n-N+1}+q^{N-1}\big)\widetilde{Q}_j, \qquad j \in \mathbb Z_+.
\end{gather*}

Thus we obtain that $P_j$ and $\widetilde{Q}_j$ are related by the change of variable
\[
\lambda=\frac{2q^{m-n+1}}{(1-q^2)^2}z-\frac{q^{m-n+1}}{(1-q^2)^2}\big(q^{N-1}+q^{-(N-1)}\big).
\]
So,
\[
P_j(\lambda)=\sqrt{\frac{1}{(q^2,q^{2n};q^2)_j}}  Q_j\big(z;q^{2n-N+1},q^{N-1}|q^2\big),
\]
where $z=\cos \theta=\frac{e^{i\theta}+e^{-i\theta}}{2}$ and
\begin{gather*}
\lambda=-\frac{q^{2-2n}}{(1-q^2)^2}\big(1-q^{N-1}e^{i\theta}\big)\big(1-q^{N-1}e^{-i\theta}\big),\\
Uf_j=q^{-j(N-1)}\sqrt{\frac{(q^{2j+2};q^2)_{n-1}}{(q^{2};q^2)_{n-1}}}\sqrt{\frac{1}{(q^2,q^{2n};q^2)_j}}
Q_j\big(z;q^{2n-N+1},q^{N-1}|q^2\big) \\
\phantom{Uf_j}{} =q^{-j(N-1)}\frac{1}{(q^2;q^2)_j} Q_j\big(z;q^{2n-N+1},q^{N-1}|q^2\big).
\end{gather*}
On the other hand,
\[
\int_1^\infty
f_j(x)\Phi_{l}(x)\rho(x)d_{q^{-2}}x=\frac{q^{-j(N-1)}}{(q^2;q^2)_j} Q_j\big(z;q^{2n-N+1},q^{N-1}|q^2\big),
\]
where $z=\frac12\big(q^{2l+N-1}+q^{-(2l+N-1)}\big)$.

Hence for every function $f(x)$ on $q^{-2\mathbb{Z}_+}$ with f\/inite support one has
\[
Uf=\int_1^\infty f(x)\Phi_{l}(x)\rho(x)d_{q^{-2}}x.
\]
Now the claim of the theorem follows from the orthogonality relations for the Al-Salam--Chihara polynomials~\eqref{rel_orth_ASC}. \end{proof}
\begin{thebibliography}{99}
\bibitem[1]{AG}
Akhiezer N.I., Glazman I.M.,
Theory of linear operators in Hilbert space,  Dover Publications, New York, 1993.

\bibitem[3]{BS}
Bershtein O., Sinelshchikov S.,
Function theory on a $q$-analog of complex hyperbolic space,
\href{http://arxiv.org/abs/1009.6063}{arXiv:1009.6063}.

\bibitem[4]{Ch-P}
Chari V., Pressley A.,
A guide to quantum groups,
Cambridge University Press, Cambridge, 1994.


\bibitem[5]{Faraut}
Faraut J.,
Distributions sph\'eriques sur les espaces hyperboliques,
{\it J.~Math. Pures Appl.} {\bf 58} (1979), 369--444.

\bibitem[7]{Jant}
Jantzen J.C.,
Lectures on quantum groups,
{\it Graduate Studies in Mathematics}, Vol.~6, American Mathema\-ti\-cal Society, Providence, RI, 1996.

\bibitem[8]{Kl-Sch}
Klimyk A., Schm\"udgen K.,
Quantum groups and their representations,
{\it Texts and Monographs in Physics}, Springer-Verlag, Berlin, 1997.

\bibitem[9]{koekoek}
Koekoek R., Swarttouw R.F.,
The Askey-scheme of hypergeometric orthogonal polynomials and its $q$-analogue, Report 98-17,
Faculty of Technical Mathematics and Informatics, Delft University of Technology, 1998, \url{http://aw.twi.tudelft.nl/~koekoek/askey/}.

\bibitem[12]{Molch2}
Molchanov V.F.,
Harmonic analysis on homogeneous spaces,
{\it Itogi Nauki i Tekhniki}, Vol.~59, VINITI, Moscow, 1990, 5--144
(English transl.:  {\it Encycl. Math. Sci.}, Vol.~59, Springer, Berlin, 1995, 1--135).

\bibitem[13]{VO}
Onishchik A.L., Vinberg E.B.,
Lie groups and algebraic groups,
{\it Springer Series in Soviet Mathematics}, Springer-Verlag, Berlin, 1990.

\bibitem[14]{RTF}
Reshetikhin N.Yu., Takhtadzhyan L.A., Faddeev L.D.,
Quantization of Lie groups and Lie algebras,
{\it Algebra i Analiz} {\bf  1} (1989), 178--206 (English transl.: {\it Leningrad Math.~J.} {\bf 1} (1990), 193--225).

\bibitem[15]{polmat}
Shklyarov D., Sinel'shchikov S., Vaksman L.,
Fock representations and quantum matrices,
\href{http://dx.doi.org/10.1142/S0129167X04002600}{{\it Internat.~J. Math.}} {\bf 15} (2004), 855--894,
\href{http://arxiv.org/abs/math.QA/0410605}{math.QA/0410605}.

\end{thebibliography}
\end{document}
